%=======================================================================
\needspace{5\baselineskip}
\section{Lessons from the Literature}
\label{sec:lit}
%=======================================================================

Carnein et al.~\cite{carnein2017comparison} tune each algorithm per dataset, offline, with irace on ARI. Mansalis et al.~\cite{mansalis2018evaluation} vary one parameter at a time and read the plots, with one run per setting. confStream~\cite{carnein2020confstream} tunes online with irace-inspired sampling but no racing, and IWFO~\cite{thakral2024iwfo} optimises batch centroids rather than hyperparameters. Four lessons carry over.
\begin{enumerate}[label=\textbf{L\arabic*.},leftmargin=2em,itemsep=1pt,topsep=2pt]
  \item \textbf{Relative values transfer better than absolute ones.} Carnein's tuned DBSTREAM radius is $0.02$ on Sensor but $1.95$ on KDDCup99 even on standardised data, while Mansalis's 15--20 micro-clusters per class is relative, like our $c\cdot s(D)$ and $q/k$. We keep the relative grid.
  \item \textbf{The metric and the ranges are design choices.} Carnein caps micro-clusters because ``the ARI can be artificially improved'' by adding more, and in Mansalis purity and CMM rank the algorithms in opposite orders. Our score and thresholds need an explicit basis (P2).
  \item \textbf{$\beta$ is a dead parameter.} MOA and River both use it only through the product $\beta\mu$, and pinning it at $0.75$ costs only $0.0004$ ARI of ceiling. We therefore pin it (P1).
  \item \textbf{Racing suits offline tuning.} confStream leaves it out because it ``makes it difficult to apply irace for streaming applications''. Our tuning is offline, so we race (P1, P2).
\end{enumerate}

\subsection{Racing and headroom}
\label{sec:thresholds}

irace~\cite{lopezibanez2016irace} evaluates candidates block by block and, after each block, drops those a statistical test finds worse than the best, so extra seeds go only to survivors. Following its guide~\cite{iraceguide2025}, a block holds one seed over all tunable datasets, and we use the t-test because the default Friedman test ranks within instances and ignores any rescaling of scores~\cite{demsar2006statistical}. The race covers exactly the grid's candidates and minimises $-h$, where
\[
  h_m(d) = \frac{\mathrm{ARI}_m(d) - \mathrm{ARI}_{\text{default}}(d)}{\mathrm{ARI}_{\text{best}}(d) - \mathrm{ARI}_{\text{default}}(d)}
\]
is the headroom proposed by Probst et al.~\cite{probst2019tunability}: $1$ matches the best grid point and $0$ is no better than the default. It addresses the concern of van Rijn and Hutter~\cite{vanrijn2018importance} that performance scales differ between datasets. Table~\ref{tab:thresholds} lists its three thresholds; the gain gate uses the detectable gain
$(z_{0.95}+z_{0.80})\,\mathrm{SE} \approx 2.49\,\mathrm{SE}$, with $\mathrm{SE} = \sqrt{(\mathrm{sd}_{\text{best}}^2+\mathrm{sd}_{\text{def}}^2)/n}$.

\begin{table}[H]
\centering
\caption{The three thresholds around headroom.}
\label{tab:thresholds}
\small
\begin{tabularx}{\textwidth}{@{}l l c L@{}}
\toprule
\textbf{Threshold} & \textbf{Purpose} & \textbf{Now} & \textbf{Basis} \\
\midrule
Gain gate & Gain above noise & $0.02$ & Statistical: $\text{gain}/\mathrm{SE} \ge z_{0.95}$ on tuning seeds \\
Ceiling gate & Gain large enough to matter & $0.05$ & Judgement, fixed before seeing the data~\cite{lakens2018equivalence} \\
Tie band & Method differs from default & 1 sd & Statistical: $1.645\sqrt{2/n}\,\mathrm{sd} = 1.04\,\mathrm{sd}$ at $n=5$ \\
\bottomrule
\end{tabularx}
\end{table}
\takeaway{Only the ceiling gate is a matter of judgement. Covertype's gain of $0.019$ is statistically real ($z \approx 18$) but has no practical meaning, because class labels need not match the cluster structure~\cite{faerber2010labels}. The inflated ceiling in Table~\ref{tab:noise} is the selection bias described by Cawley and Talbot~\cite{cawley2010overfitting}.}
